\documentclass[11pt]{article}

\usepackage[margin=1in]{geometry}
\usepackage[utf8]{inputenc}
\usepackage[T1]{fontenc}
\usepackage{lmodern}
\usepackage{amsmath,amssymb,amsthm}
\usepackage{booktabs}
\usepackage{microtype}
\usepackage[numbers,sort&compress]{natbib}
\usepackage{xcolor}
\usepackage[colorlinks=true,linkcolor=blue!50!black,citecolor=blue!50!black,urlcolor=blue!50!black]{hyperref}

\newcommand{\draftnote}[1]{\textcolor{red!70!black}{[\textbf{Draft:} #1]}}

\title{Freezing, Restarts and the Batch:\\What the Literature Says About Crystallising a Snapped State\\[0.5em]
\large Research note, 22 October 2026: not for circulation}

\author{
  David Barkhuizen\\
  Independent researcher
  \and
  Claude (Opus 5.5)\\
  Anthropic
}

\date{}

\begin{document}

\maketitle

\draftnote{A stand-alone note written at the first author's request, after a literature review of the
concepts and findings that arose in the continuous crystallisation between 20 and 22 October 2026
(its note is \texttt{continuous-crystallisation.tex}). The review itself, paper by paper, is \texttt{indrajala-ml}'s
\texttt{docs/\allowbreak window-recurrent-\allowbreak literature.md} (its extension of 22 October); its papers are in the
literature repository. The findings it maps come from \texttt{docs/\allowbreak state-formation-\allowbreak workplan.md}
(decisions D48--D66). Section~\ref{sec:restarts}'s measurement is new, offline, on completed runs.
Several registered probes were running at its writing; none of their data is used here.}

\begin{abstract}
Our method learns a finite-state transducer from input--output examples with a recurrent state
snapped to learned codes, cooled by a quasi-static profile, and certifies the learned machine
exactly. Every run that fails ends the same way: the profile's clock holds for the rest of the
budget, a freeze. We read the literature on annealing schedules, restarts, training dynamics and the
learnability of random automata against that finding and the others of the last two days. Three
independent results say no cooling schedule passes a freeze: the dissipation bound of finite-time
thermodynamics grows with the relaxation time, the optimal (constant-speed) schedule halts as it
diverges, and its derivation assumes a uniform system, which a glass is not. Our clock is a
bang-bang form of that optimal schedule. Restart theory, by contrast, applies exactly: a certified
run is a Las Vegas run, so the expected cost per certified machine of fresh-seed restarts at a fixed
cutoff, and its optimal cutoff, follow from the runs' own times. Measured on the completed runs,
restarting at the optimal cutoff would cost 4.4 times less per certificate than running every seed out
on $S_3$, 1.9 times on move-to-front-4 and about 10 times on $S_4$ (small samples). Annealing an
ensemble of copies at one temperature, a precedent in the literature, gives our statistical method,
a batch with a clear majority certifying, a physical reading. Neither the collapses at a steady
coupling nor the random machines' failures are explained by the work found.
\end{abstract}

\section{Setting}

The method and its terms are those of \texttt{continuous-crystallisation.tex}. A window-recurrent
network carries a factored state, each factor's soft codes annealed from a melt: a temperature per
factor, a coupling from the continuous state to the codes read from the codes' sharpness, and one
clock that advances only while the answers' screen keeps up with the level it held in the hold
(D49). A run stops when its snapped machine answers right, waits until no near tie splits it (the
margin stop, D59), and succeeds when it is certified at both verdicts (D41). The method is
statistical: a parallel batch of fresh seeds, a clear majority certifying, not every seed.

The findings this note reads against the literature (all offline on completed runs,
22 October):
\begin{itemize}
\item \textbf{The freeze.} All 16 spent runs of the continuous probes hold the clock at zero for
738--1{,}472 epochs. The random machines freeze before the state forms, near the factors' melting
point (0.54--1.16 of it); the coupling keeps rising while the clock is held. One run thawed after
445 epochs and certified. A hold in the window between melt and freeze (two runs) did not form.
\item \textbf{The collapse.} On $S_4$ the state formed (strict-exact at steps 130--134), then the
screen fell from 0.99 to 0.42 in one or two epochs at a steady coupling of 0.97, both factors
together, and the run froze.
\item \textbf{Forming against keeping.} The random machines fail at forming, $S_4$ at keeping; the
machine's size is not the variable (move-to-front-4, 24 states, is the best backdrop).
\item \textbf{The factors fail together.} In nine frozen runs and both $S_4$ collapses, no run formed
one factor while another stayed stuck; only their live readings diverge.
\end{itemize}

\section{The freeze is the known end of an equilibrium-gated schedule}

Simulated annealing named the phenomenon: cooled in slow stages, a system ``freezes'' and no further
changes occur; cooled fast, it is frozen into a metastable state, and the best of several quenches is
reported beside the anneal \citep{kirkpatrick-1983-optimization-by-simulated-annealing}. Adaptive
schedules cool only as fast as the system relaxes
\citep{huang-1986-efficient-cooling-schedule,romeo-1989-simulated-annealing-theory-layout}, and
finite-time thermodynamics made that precise. The availability a near-equilibrium process dissipates
is at least $L^2 \bar\varepsilon / \tau$: its thermodynamic length squared, times a mean relaxation
time, over the total time \citep{salamon-1983-thermodynamic-length-dissipated-availability}. Holding
the system's distribution at a constant distinguishability from the equilibrium distribution at the
current temperature gives the schedule of constant thermodynamic speed,
\begin{equation}
\frac{dT}{dt} = -\frac{vT}{\varepsilon\sqrt{C}},
\end{equation}
with heat capacity $C$ and relaxation time $\varepsilon$ estimated from the run
\citep{nulton-1988-statistical-mechanics-combinatorial-optimization,salamon-1988-simulated-annealing-constant-thermodynamic-speed};
it is derived from first principles as the leading term of the entropy-minimising path, for uniform
systems whose intensive quantities are uniquely defined
\citep{andresen-1993-constant-thermodynamic-speed-cooling,andresen-1994-constant-thermodynamic-speed-entropy-production},
and compared against exponential, logarithmic and linear cooling it is the best
\citep{nourani-1998-comparison-annealing-cooling-strategies}.

Three consequences bear on the freeze. As $\varepsilon$ diverges, which is what a glass is, the bound
diverges and the optimal schedule halts; the derivation's uniformity fails there too. So a slower
or better-shaped cool is not a remedy, which is what the order-biased cool's stage~0 found (D65,
set aside: the failures freeze; the cool does not outrun them). Second, our clock is a degenerate
form of the optimal schedule: it reads the same lag (the screen against the hold's level) but stops
at a set lag, where the optimum moves at one. Both stop at a glass; a constant-lag clock is the
principled form should the cool's pace matter again. Third, under-parametrised networks show typical
glassy dynamics \citep{baity-jesi-2018-dynamics-dnn-versus-glassy-systems}, and long gradient training
on changing data loses plasticity, with dead units, growing weights and falling effective rank as
its correlates and a non-gradient component as its remedy
\citep{dohare-2023-maintaining-plasticity-deep-continual-learning,ash-2020-warm-starting-neural-network-training}.
Our curriculum's growing length is such a change; whether our frozen runs show those correlates is
an offline reading not yet made.

What the literature does not have is the coupling quenching at a held clock: a glass forming at a
fixed temperature. That is a property of reading the coupling from the codes' sharpness (D50), not
of annealing.

\section{Restarts: the cost per certificate}
\label{sec:restarts}

The field's other answers to a frozen run are to re-anneal it, restart it, or exchange it with a
hotter copy. Re-annealing resets each parameter's annealing time from its current sensitivity
\citep{ingber-1989-very-fast-simulated-reannealing}; warm restarts with doubling periods are standard
in gradient training \citep{loshchilov-2017-sgdr-warm-restarts}, which is the slowed re-cool's
schedule (D66, a forward rate times $2^{-k}$ after the $k$-th reheat), triggered here by the freeze
rather than a fixed period. A per-factor re-anneal was considered and dropped: the factors fail
together.

Restart theory applies to our method exactly. A certified run is a Las Vegas run: whenever it stops
with a certificate, its machine is proved. For runs whose time to success $T$ has distribution $q$,
the optimal strategy is a fixed cutoff $t^*$ repeated with fresh seeds, minimising the expected cost
per success
\begin{equation}
\ell(t) = \frac{\mathbb{E}[\min(T, t)]}{q(t)},
\end{equation}
and it stays optimal when runs may be suspended and resumed
\citep{luby-1993-optimal-speedup-las-vegas-algorithms}. Run times of randomised search are often
heavy-tailed, with a left tail of short successes that rapid restarts exploit
\citep{gomes-2000-heavy-tailed-phenomena-satisfiability-csp}. So kill-and-resample (D63~(e)) needs no
fixed window: $t^*$ comes from the runs. And the comparison D63 plans, the reheat against
resampling, is $\ell$ against $\ell$: the reheat is a different algorithm with its own $q$.

Table~\ref{tab:restarts} reads $\ell$ on the completed runs of the current method (D54's cell with the
margin stop: D54's and D55's seeds, D59's re-runs where a run stopped strict-only), on the step clock,
which is the same on every machine. A run spent at its budget is censored, and $\ell$ is read only to
the shortest censored run's step.

\begin{table}[t]
\centering
\caption{Fresh-seed restarts at the optimal fixed cutoff against running every seed out: Luby's
$\ell(t^*)$ in steps per certificate, offline on completed runs (\texttt{modular\_family\_probe.py
restart-cost}). Small samples: $t^*$ is an estimate at the edge of the observed successes.}
\label{tab:restarts}
\begin{tabular}{@{}lrrrrrrr@{}}
\toprule
backdrop & runs & certified & censored & $t^*$ & $q(t^*)$ & $\ell(t^*)$ & each run out \\
\midrule
$S_3$ & 24 & 17 & 7 & 137 & 0.62 & 182 & 798 \\
move-to-front-4 & 11 & 10 & 1 & 169 & 0.91 & 157 & 298 \\
$S_4$ & 3 & 1 & 2 & 106 & 0.33 & 318 & 3{,}213 \\
\bottomrule
\end{tabular}
\end{table}

The successes cluster between about 85 and 170 steps; $\ell$ is least just after the cluster and rises
slowly beyond it ($S_3$: 182 at step 137, 216 at 219). The one late success ($S_3$ seed 30, certified
at step 854 after a 445-epoch freeze) costs more to wait for than it returns. A cutoff of 200 steps,
the first sizing, costs about 15--20\% more than $S_3$'s estimated optimum and guards the slower
successes. The random backdrops (0 of 4 and 1 of 4) are too thin for a curve. Luby's model counts only
the time to the first success: a certificate that only a reheat reaches, after a collapse and a
recovery, is outside it and is read apart, as D63 plans.

\section{The batch as an ensemble}

Constant-speed annealing was implemented on an ensemble of $n$ independent copies sharing one
temperature, the lag measured across them, because a single run is not ergodic on annealing's time
scales \citep{salamon-1988-simulated-annealing-constant-thermodynamic-speed}. Two of its remarks bear
on us. The ensemble stops representing equilibrium once phase space breaks into more basins than
there are copies, expected for glassy problems: a mechanism for the scaling conjecture's remedy for
random machines, a larger batch. And on problems without local minima, a direct quench beats any
equilibrium-following schedule: ordered machines might certify under much faster cools. Exchanging
configurations between copies at a ladder of temperatures, parallel tempering, is the batch-level
form of escape \citep{earl-2005-parallel-tempering-theory-applications}; its exactness (detailed
balance) does not carry to gradient training, which is no Metropolis chain.

\section{Collapses, smoothing and random machines}

\textbf{Collapses.} With adaptive optimisers, after a network fits its data, the last layer's norm
grows and the loss spikes, in cycles \citep{thilak-2022-slingshot-mechanism-adaptive-optimizers}; full
batch gradient descent hovers at the edge of stability with non-monotone spikes
\citep{cohen-2021-edge-of-stability-gradient-descent} (their ``sharpness'', the Hessian's top
eigenvalue, is not ours). We train with Adam, and $S_4$'s collapse fits the slingshot's timing; it is
untested, because the readout's weight norm is not logged. If it is an optimiser's instability, the
paced coupling (D61), which smooths the coupling's reading, cannot cure it.

\textbf{Smoothing.} The Savitzky--Golay filter is a local least-squares polynomial read at the
window's centre \citep{savitzky-1964-smoothing-differentiation-least-squares}; centred, degree one is a
moving average, and the cutoff frequency follows from the degree and the window
\citep{schafer-2011-what-is-savitzky-golay-filter}. D61's smoothed arm reads a degree-one line at the
newest epoch, which neither treats: its offline lag advantage over a moving average comes from that
one-sided reading, a linear extrapolation, whose variance at the window's end is about four times
that at its centre. The arm may make coupling moves of its own from noise; its reading should check.

\textbf{Random machines.} Random automata are hard for statistical queries under adversarial input
distributions \citep{angluin-2010-lower-bounds-random-structures-sq}, and random semiautomata under
the uniform distribution once words reach $\Omega(N^2 \ln N)$ symbols with only the final state
labelled \citep{giapitzakis-2025-sq-complexity-learning-semiautomata}; coarse gradient training is no
stronger than statistical queries \citep{abbe-2021-power-differentiable-learning-pac-sq}. With state
labels, random automata are easy \citep{angluin-2015-learning-random-dfa-uniform-strings-state}. Our
words are under 20 symbols with an answer at every step, outside every regime found: the random
machines' failure to form is not explained by this theory.

\section{Reading the results}

The statistics our claims need are those the papers already use: Wilson's interval for a rate
\citep{wilson-1927-probable-inference}, Fisher's exact test for an arm against its control
\citep{fisher-1922-contingency}, Holm's procedure across backdrops
\citep{holm-1979-sequentially-rejective}, and Mann and Whitney's test for costs
\citep{mann-whitney-1947}. The scaling conjecture's claim, that the rate falls as machines grow, is a
trend in proportions across ordered backdrops, which Cochran and Armitage's test reads
\citep{cochran-1954-chi-square,armitage-1955-trend}. Two gaps: D63's reheat and resampling share
seeds, a paired comparison (McNemar's test, not yet in the literature repository), and run times
censored at the budget suit a log-rank test better than Mann and Whitney's. Seeds are a source of
variance to randomise, compared by their distribution, not their means
\citep{bouthillier-2021-accounting-variance-ml-benchmarks}.

\section{Direction}

\begin{enumerate}
\item \textbf{The schedule is not the lever.} The cool stays as it is; failures are escaped (a reheat
or a fresh seed) and the method is judged by its cost per certified machine.
\item \textbf{The batch is the unit.} Fresh-seed restarts at $t^*$, read per backdrop from the runs;
larger batches where basins are many (the random machines).
\item \textbf{Offline, now:} $\ell$ on every backdrop as runs complete (done for three); plasticity's
correlates on the frozen runs; when D61 is read, the smoothed arm's noise; when D64 is read, a trend
test of the scaling conjecture.
\item \textbf{Measurements to build into the next change of the per-epoch row:} the readout's weight
norm, for the slingshot.
\item \textbf{Later candidates, not planned:} a constant-lag clock; fast cools for ordered machines;
parallel tempering across a batch.
\end{enumerate}

\section{What is not established}

None of the literature's mechanisms is shown to be ours. The freeze is consistent with a kinetic
arrest, but no relaxation time was measured diverging; the slingshot reading of $S_4$'s collapse and
the plasticity reading of the freeze are untested; $t^*$ rests on 24, 11 and 3 runs, and the restart
model assumes runs independent and identically distributed, which fresh seeds are by construction
but different machines' clocks are not. The search for literature cannot show absence: work under
other names may have been missed.

\section{Provenance}

The review: \texttt{indrajala-ml} \texttt{docs/\allowbreak window-recurrent-\allowbreak literature.md}, its extension of
22 October (\#928--\#930, \#932); the papers in the literature repository (\#8--\#19). The findings
mapped: \texttt{docs/\allowbreak state-formation-\allowbreak workplan.md}, ``the scaling conjecture'' and ``the freeze''
(\#922, \#923), D61--D66, and their update from this review (\#934). The restart cost: \texttt{modular\_family\_probe.py restart-cost} (\#933) on the
completed logs of D54, D55, D57 and D59 (the move-to-front-4 seed of D54 still running left out).

\bibliographystyle{plainnat}
\bibliography{library,extra}

\end{document}
